\documentclass[aspectratio=169,10pt]{beamer}
% ═══════════════════════════════════════════════════════════════════════════════
% SHARED PREAMBLE — Foundation Models Course (HIT)
% To be \input by each lecture file AFTER \documentclass
% ═══════════════════════════════════════════════════════════════════════════════

% ─────────────────────────────────────────────────────────────────────────────
% BEAMER THEME & BASIC SETUP
% ─────────────────────────────────────────────────────────────────────────────
\usetheme{Madrid}
\usecolortheme{default}

% Remove navigation symbols
\beamertemplatenavigationsymbolsempty

% ─────────────────────────────────────────────────────────────────────────────
% COLORS — HIT Branding
% ─────────────────────────────────────────────────────────────────────────────
\definecolor{hitblue}{RGB}{0,51,102}        % Primary: HIT Navy
\definecolor{hitgold}{RGB}{192,148,0}       % Accent: HIT Gold
\definecolor{hitlight}{RGB}{230,240,250}    % Light background

% Override Beamer palette
\setbeamercolor{primary}{fg=white,bg=hitblue}
\setbeamercolor{secondary}{fg=hitblue,bg=hitlight}
\setbeamercolor{structure}{fg=hitblue}
\setbeamercolor{alerted text}{fg=hitgold}

% ─────────────────────────────────────────────────────────────────────────────
% PACKAGES
% ─────────────────────────────────────────────────────────────────────────────
\usepackage{amsmath}
\usepackage{amssymb}
\usepackage{mathtools}
\usepackage{booktabs}
\usepackage{graphicx}
\usepackage{hyperref}
\usepackage{tikz}
\usepackage{pgfplots}
\usepackage{tcolorbox}
\usepackage{listings}
\usepackage{xcolor}
\usepackage{algorithm}
\usepackage{algpseudocode}

% ─────────────────────────────────────────────────────────────────────────────
% TikZ LIBRARIES
% ─────────────────────────────────────────────────────────────────────────────
\usetikzlibrary{shapes,arrows,positioning,calc,chains,scopes}
\pgfplotsset{compat=1.17}

% ─────────────────────────────────────────────────────────────────────────────
% CODE LISTINGS
% ─────────────────────────────────────────────────────────────────────────────
\lstset{
  basicstyle=\ttfamily\small,
  breaklines=true,
  commentstyle=\color{gray},
  keywordstyle=\color{hitblue},
  stringstyle=\color{hitgold},
  showstringspaces=false,
  backgroundcolor=\color{hitlight},
  frame=single,
  rulecolor=\color{hitblue}
}

% ─────────────────────────────────────────────────────────────────────────────
% TCOLORBOX STYLES
% ─────────────────────────────────────────────────────────────────────────────
\tcbset{
  hitbox/.style={
    colback=hitlight,
    colframe=hitblue,
    fonttitle=\bfseries,
    boxrule=1pt
  },
  alertbox/.style={
    colback=yellow!10,
    colframe=hitgold,
    fonttitle=\bfseries,
    boxrule=1.5pt
  }
}

% ─────────────────────────────────────────────────────────────────────────────
% CUSTOM COMMANDS
% ─────────────────────────────────────────────────────────────────────────────

% Placeholder for figures
\newcommand{\placeholder}[3]{
  \begin{center}
    \fbox{\begin{minipage}[c][#2][c]{#1}
      \centering\small\color{gray}[Figure: #3]
    \end{minipage}}
  \end{center}
}

% Section divider frame
\newcommand{\sectionframe}[2]{
  \begin{frame}[plain]
    \begin{tikzpicture}[remember picture,overlay]
      \fill[hitblue] (current page.south west) rectangle (current page.north east);
      \node[anchor=center, white, font=\Large\bfseries] at (current page.center) {
        \begin{tabular}{c}
          #1 \\[0.3cm]
          {\small\color{hitgold}#2}
        \end{tabular}
      };
    \end{tikzpicture}
  \end{frame}
}

% ─────────────────────────────────────────────────────────────────────────────
% LOAD CUSTOM MACROS
% ─────────────────────────────────────────────────────────────────────────────
% ═══════════════════════════════════════════════════════════════════════════════
% SHARED MACROS — Mathematical Notation for Foundation Models Course
% ═══════════════════════════════════════════════════════════════════════════════

% ─────────────────────────────────────────────────────────────────────────────
% ACTIVATION & LAYER FUNCTIONS
% ─────────────────────────────────────────────────────────────────────────────
\newcommand{\softmax}{\mathrm{softmax}}
\newcommand{\sigmoid}{\mathrm{sigmoid}}
\newcommand{\relu}{\mathrm{ReLU}}
\newcommand{\gelu}{\mathrm{GELU}}
\newcommand{\LayerNorm}{\mathrm{LayerNorm}}
\newcommand{\FFN}{\mathrm{FFN}}
\newcommand{\MHA}{\mathrm{MHA}}
\newcommand{\Attn}{\mathrm{Attention}}

% ─────────────────────────────────────────────────────────────────────────────
% ATTENTION COMPONENTS
% ─────────────────────────────────────────────────────────────────────────────
\newcommand{\query}{Q}
\newcommand{\key}{K}
\newcommand{\val}{V}
\newcommand{\dmodel}{d_{\text{model}}}
\newcommand{\dk}{d_k}
\newcommand{\nheads}{h}

% Scaled dot-product attention formula
\newcommand{\SDPA}{
  \Attn(\query, \key, \val) = \softmax\left(\frac{\query \key^T}{\sqrt{\dk}}\right) \val
}

\newcommand{\CrossAttn}{\text{CrossAttention}}

% ─────────────────────────────────────────────────────────────────────────────
% PROBABILITY & EXPECTATION
% ─────────────────────────────────────────────────────────────────────────────
\newcommand{\E}{\mathbb{E}}
\newcommand{\Prob}{\mathbb{P}}
\newcommand{\KL}{\mathrm{KL}}
\newcommand{\ELBO}{\mathrm{ELBO}}
\newcommand{\given}{\mid}

% ─────────────────────────────────────────────────────────────────────────────
% NORMS & OPERATIONS
% ─────────────────────────────────────────────────────────────────────────────
\newcommand{\norm}[1]{\left\lVert #1 \right\rVert}
\newcommand{\abs}[1]{\left| #1 \right|}
\newcommand{\inner}[2]{\langle #1, #2 \rangle}
\newcommand{\T}{^{\top}}

% ─────────────────────────────────────────────────────────────────────────────
% VECTORS & MATRICES
% ─────────────────────────────────────────────────────────────────────────────
\newcommand{\bvec}[1]{\mathbf{#1}}
\newcommand{\mat}[1]{\mathbf{#1}}
\newcommand{\iid}{\stackrel{\text{i.i.d.}}{\sim}}
\newcommand{\defeq}{\coloneq}

% ─────────────────────────────────────────────────────────────────────────────
% LANGUAGE MODELING
% ─────────────────────────────────────────────────────────────────────────────
\newcommand{\vocab}{\mathcal{V}}
\newcommand{\context}{\text{context}}
\newcommand{\token}{t}
\newcommand{\seq}{x}
\newcommand{\ptheta}{p_\theta}
\newcommand{\pref}{p_{\text{ref}}}

% ─────────────────────────────────────────────────────────────────────────────
% RLHF & DPO
% ─────────────────────────────────────────────────────────────────────────────
\newcommand{\piref}{\pi_{\text{ref}}}
\newcommand{\pitheta}{\pi_\theta}
\newcommand{\yw}{y_w}
\newcommand{\yl}{y_l}
\newcommand{\DPOloss}{\mathcal{L}_{\mathrm{DPO}}}

% ─────────────────────────────────────────────────────────────────────────────
% RETRIEVAL & RAG
% ─────────────────────────────────────────────────────────────────────────────
\newcommand{\docs}{\mathcal{D}}
\newcommand{\qvec}{q}
\newcommand{\dvec}{d}

% ─────────────────────────────────────────────────────────────────────────────
% AGENTS & REASONING
% ─────────────────────────────────────────────────────────────────────────────
\newcommand{\obs}{o}
\newcommand{\act}{a}
\newcommand{\thought}{t}
\newcommand{\tools}{\mathcal{T}}
\newcommand{\history}{h}
\newcommand{\concat}{\oplus}

% ─────────────────────────────────────────────────────────────────────────────
% SETS & LOGIC
% ─────────────────────────────────────────────────────────────────────────────
\newcommand{\R}{\mathbb{R}}
\newcommand{\N}{\mathbb{N}}
\newcommand{\Z}{\mathbb{Z}}

\endinput


% ─────────────────────────────────────────────────────────────────────────────
% TITLE & AUTHOR (can be overridden in individual lectures)
% ─────────────────────────────────────────────────────────────────────────────
\author[D. Lederman]{Dror Lederman}
\institute{Holon Institute of Technology (HIT) \\ Data Science Department}
\date{\today}

% ─────────────────────────────────────────────────────────────────────────────
% FOOTER & HEADER
% ─────────────────────────────────────────────────────────────────────────────
\setbeamertemplate{footline}{
  \leavevmode%
  \hbox{%
    \begin{beamercolorbox}[wd=0.33\paperwidth,ht=2.25ex,dp=1ex,center]{author in head/foot}%
      \usebeamerfont{author in head/foot}\insertshortauthor
    \end{beamercolorbox}%
    \begin{beamercolorbox}[wd=0.34\paperwidth,ht=2.25ex,dp=1ex,center]{title in head/foot}%
      \usebeamerfont{title in head/foot}\insertshorttitle
    \end{beamercolorbox}%
    \begin{beamercolorbox}[wd=0.33\paperwidth,ht=2.25ex,dp=1ex,right]{frame number}%
      \usebeamerfont{frame number}\insertframenumber{} / \inserttotalframenumber\hspace*{2ex}
    \end{beamercolorbox}%
  }%
  \vskip0pt%
}

\endinput


\title{Week 3: Training Large Language Models}
\subtitle{Foundation Models and Agentic AI}
\author{Lecturer Name}
\institute{Holon Institute of Technology (HIT)}
\date{Semester, 2025}

\begin{document}

\begin{frame}
  \titlepage
\end{frame}

\begin{frame}{Learning Objectives}
  After this lecture you will be able to:
  \begin{itemize}
    \item Compare \textbf{pretraining objectives}: causal LM, masked LM, span corruption.
    \item Describe \textbf{data curation} practices for large-scale pretraining.
    \item Explain \textbf{supervised fine-tuning (SFT)} and instruction tuning.
    \item Derive the \textbf{RLHF} pipeline: reward model training and PPO optimization.
    \item Understand \textbf{DPO} and how it eliminates the explicit reward model.
  \end{itemize}
\end{frame}

\section{Pretraining Objectives}
\sectionframe{Pretraining Objectives}{What does the model learn to predict?}

\begin{frame}{Causal Language Modeling (CLM)}
  \textbf{Used by:} GPT family, Llama, Claude, Mistral.
  \begin{tcolorbox}[hitbox, title=Objective: Next-Token Prediction]
    \[
      \mathcal{L}_\text{CLM}(\theta) = -\sum_{t=1}^{T} \log \ptheta(x_t \given x_1, \ldots, x_{t-1})
    \]
  \end{tcolorbox}
  \vspace{0.4em}
  \begin{itemize}
    \item Training signal: predict the \emph{next} token given all previous tokens.
    \item Uses a \textbf{causal mask}: token $t$ can only attend to tokens $1, \ldots, t-1$.
    \item Simple yet powerful: implicitly learns syntax, semantics, and world knowledge.
    \item Naturally supports \textbf{autoregressive generation}.
  \end{itemize}
  \note{This is the workhorse of modern LLMs. Ask students: why does predicting the next word require understanding language at a deep level?}
\end{frame}

\begin{frame}{Masked Language Modeling (MLM)}
  \textbf{Used by:} BERT, RoBERTa, DeBERTa.
  \begin{tcolorbox}[hitbox, title=Objective: Fill in the Blank]
    Randomly mask 15\% of tokens; predict the original tokens:
    \[
      \mathcal{L}_\text{MLM}(\theta) = -\sum_{i \in \mathcal{M}} \log \ptheta(x_i \given \bvec{x}_{\setminus \mathcal{M}})
    \]
  \end{tcolorbox}
  \vspace{0.4em}
  \textbf{Pros:} Bidirectional context → better representations for classification/NER.\\
  \textbf{Cons:} Cannot generate autoregressively; train-test discrepancy (no \texttt{[MASK]} at inference).
  \vspace{0.4em}\\
  \textbf{T5 variant:} Span corruption — mask \emph{contiguous spans}, not individual tokens.
\end{frame}

\section{Data Curation}
\sectionframe{Data Curation}{Garbage in, garbage out — at trillion-token scale}

\begin{frame}{Pretraining Data Pipeline}
  \begin{columns}[T]
    \column{0.52\textwidth}
      \textbf{Raw data sources:}
      \begin{itemize}
        \item Common Crawl (trillions of tokens, noisy)
        \item Books (BookCorpus, Project Gutenberg)
        \item Wikipedia, arXiv, GitHub
        \item Curated web (C4, The Pile, RedPajama)
      \end{itemize}
      \vspace{0.5em}
      \textbf{Filtering steps:}
      \begin{enumerate}
        \item Language identification
        \item Quality filtering (perplexity, heuristics)
        \item Near-deduplication (MinHash LSH)
        \item Safety filtering (pornography, hate speech)
        \item Domain upsampling / weighting
      \end{enumerate}
    \column{0.44\textwidth}
      \placeholder{0.95\textwidth}{4cm}{Data pipeline: web → filter → deduplicate → tokenize}
  \end{columns}
  \note{The Pile paper (Gao et al. 2020) provides detailed statistics. Llama 2's data section shows how domain weighting affects downstream tasks.}
\end{frame}

\begin{frame}{Tokenization}
  \begin{columns}[T]
    \column{0.55\textwidth}
      \textbf{Byte Pair Encoding (BPE):}
      \begin{itemize}
        \item Start: single characters as vocabulary.
        \item Iteratively merge the most frequent pair.
        \item Result: subword vocabulary ($\sim$50k tokens).
      \end{itemize}
      \vspace{0.4em}
      \textbf{SentencePiece (Unigram LM):}
      \begin{itemize}
        \item Language-independent; handles any script.
        \item Used by T5, Llama.
      \end{itemize}
      \vspace{0.4em}
      \textbf{Impact:} Tokenization affects downstream task performance, multilingual capability, and math reasoning.
    \column{0.40\textwidth}
      \placeholder{0.90\textwidth}{3.5cm}{BPE: ``foundation'' → [``found'', ``ation'']}
  \end{columns}
\end{frame}

\section{Supervised Fine-Tuning \& Instruction Tuning}
\sectionframe{Supervised Fine-Tuning}{From a language model to a helpful assistant}

\begin{frame}{Instruction Tuning / SFT}
  \textbf{Goal:} Teach the model to follow natural-language instructions.
  \vspace{0.4em}
  \begin{tcolorbox}[hitbox, title=SFT Procedure]
    \begin{enumerate}
      \item Collect demonstration data: $\{(\text{instruction}_i, \text{response}_i)\}_{i=1}^N$
      \item Fine-tune the pretrained model with CLM loss on response tokens only.
      \item Result: model that \emph{follows instructions}, not just completes text.
    \end{enumerate}
  \end{tcolorbox}
  \vspace{0.4em}
  \textbf{Examples:}
  \begin{itemize}
    \item \textit{FLAN} (Wei et al.\ 2022~\cite{wei2022finetuned}): 60+ NLP tasks → strong zero-shot generalization.
    \item \textit{InstructGPT} (Ouyang et al.\ 2022~\cite{ouyang2022training}): human-written demonstrations + RLHF.
    \item \textit{Alpaca}, \textit{Vicuna}: open-source SFT on GPT-generated instructions.
  \end{itemize}
  \note{SFT alone shifts the model's style but does not guarantee helpfulness or harmlessness. That requires RLHF.}
\end{frame}

\section{Reinforcement Learning from Human Feedback (RLHF)}
\sectionframe{RLHF}{Learning from human preferences}

\begin{frame}{Why RLHF?}
  \begin{tcolorbox}[alertbox, title=Problem with SFT Alone]
    A model trained only on demonstrations will learn to \emph{imitate} responses —
    but human preferences are often complex and hard to demonstrate.
    We need a signal that captures: ``this response is \textbf{better} than that one.''
  \end{tcolorbox}
  \vspace{0.5em}
  \textbf{RLHF} (Christiano et al.\ 2017; Ouyang et al.\ 2022):
  \begin{enumerate}
    \item Collect human \textbf{preference pairs}: $(x, y_w, y_l)$ — prompt, preferred, rejected.
    \item Train a \textbf{reward model} $r_\phi(x, y)$ on the preference data.
    \item Fine-tune the LLM with RL (PPO) to maximize $r_\phi$.
  \end{enumerate}
  \vspace{0.3em}
  % tikz/rlhf_pipeline.tex
% Three-phase RLHF pipeline: SFT → Reward Model → PPO
% Usage: % tikz/rlhf_pipeline.tex
% Three-phase RLHF pipeline: SFT → Reward Model → PPO
% Usage: % tikz/rlhf_pipeline.tex
% Three-phase RLHF pipeline: SFT → Reward Model → PPO
% Usage: \input{../../tikz/rlhf_pipeline}
\begin{tikzpicture}[
  node distance=1.0cm and 0.9cm,
  phase/.style={rectangle, draw=hitblue, fill=hitlightblue, rounded corners=4pt,
                minimum width=2.8cm, minimum height=1.6cm, font=\small, align=center,
                text width=2.6cm},
  data/.style={rectangle, draw=hitgold, fill=hitgold!12, rounded corners=3pt,
               minimum width=2.2cm, minimum height=0.55cm, font=\scriptsize, align=center},
  arr/.style={->, >=Stealth, very thick, hitblue},
  darr/.style={->, >=Stealth, thick, hitgold, dashed},
  label/.style={font=\scriptsize, gray, align=center},
]

% Phase 1: SFT
\node[phase] (sft) {\textbf{Phase 1}\\[4pt]Supervised\\Fine-Tuning\\(SFT)};

% Phase 2: Reward Model
\node[phase, right=1.5cm of sft] (rm) {\textbf{Phase 2}\\[4pt]Reward Model\\Training\\(RM)};

% Phase 3: PPO
\node[phase, right=1.5cm of rm] (ppo) {\textbf{Phase 3}\\[4pt]Reinforcement\\Learning\\(PPO)};

% Data nodes above phases
\node[data, above=0.7cm of sft] (d1) {Demonstration\\data};
\node[data, above=0.7cm of rm] (d2) {Human\\preference pairs};
\node[data, above=0.7cm of ppo] (d3) {Prompts\\(unlabeled)};

% Output nodes below phases
\node[data, below=0.7cm of sft] (o1) {$\pi_\text{SFT}$};
\node[data, below=0.7cm of rm] (o2) {$r_\phi(x,y)$};
\node[data, below=0.7cm of ppo] (o3) {$\pi_\theta$ (RLHF)};

% Arrows: data → phase
\draw[darr] (d1) -- (sft);
\draw[darr] (d2) -- (rm);
\draw[darr] (d3) -- (ppo);

% Arrows: phase → output
\draw[arr] (sft) -- (o1);
\draw[arr] (rm) -- (o2);
\draw[arr] (ppo) -- (o3);

% Arrows: phase to phase
\draw[arr] (sft) -- node[above, label] {initialize} (rm);
\draw[arr] (rm) -- node[above, label] {reward\\signal} (ppo);

% KL penalty annotation
\draw[darr, bend left=30] (o1.south) to
  node[below, font=\scriptsize, gray] {KL penalty $\beta$} (o3.south);

% Phase labels below diagram
\node[label, below=0.15cm of o1] {SFT model};
\node[label, below=0.15cm of o2] {Reward model};
\node[label, below=0.15cm of o3] {Aligned LLM};

\end{tikzpicture}

\begin{tikzpicture}[
  node distance=1.0cm and 0.9cm,
  phase/.style={rectangle, draw=hitblue, fill=hitlightblue, rounded corners=4pt,
                minimum width=2.8cm, minimum height=1.6cm, font=\small, align=center,
                text width=2.6cm},
  data/.style={rectangle, draw=hitgold, fill=hitgold!12, rounded corners=3pt,
               minimum width=2.2cm, minimum height=0.55cm, font=\scriptsize, align=center},
  arr/.style={->, >=Stealth, very thick, hitblue},
  darr/.style={->, >=Stealth, thick, hitgold, dashed},
  label/.style={font=\scriptsize, gray, align=center},
]

% Phase 1: SFT
\node[phase] (sft) {\textbf{Phase 1}\\[4pt]Supervised\\Fine-Tuning\\(SFT)};

% Phase 2: Reward Model
\node[phase, right=1.5cm of sft] (rm) {\textbf{Phase 2}\\[4pt]Reward Model\\Training\\(RM)};

% Phase 3: PPO
\node[phase, right=1.5cm of rm] (ppo) {\textbf{Phase 3}\\[4pt]Reinforcement\\Learning\\(PPO)};

% Data nodes above phases
\node[data, above=0.7cm of sft] (d1) {Demonstration\\data};
\node[data, above=0.7cm of rm] (d2) {Human\\preference pairs};
\node[data, above=0.7cm of ppo] (d3) {Prompts\\(unlabeled)};

% Output nodes below phases
\node[data, below=0.7cm of sft] (o1) {$\pi_\text{SFT}$};
\node[data, below=0.7cm of rm] (o2) {$r_\phi(x,y)$};
\node[data, below=0.7cm of ppo] (o3) {$\pi_\theta$ (RLHF)};

% Arrows: data → phase
\draw[darr] (d1) -- (sft);
\draw[darr] (d2) -- (rm);
\draw[darr] (d3) -- (ppo);

% Arrows: phase → output
\draw[arr] (sft) -- (o1);
\draw[arr] (rm) -- (o2);
\draw[arr] (ppo) -- (o3);

% Arrows: phase to phase
\draw[arr] (sft) -- node[above, label] {initialize} (rm);
\draw[arr] (rm) -- node[above, label] {reward\\signal} (ppo);

% KL penalty annotation
\draw[darr, bend left=30] (o1.south) to
  node[below, font=\scriptsize, gray] {KL penalty $\beta$} (o3.south);

% Phase labels below diagram
\node[label, below=0.15cm of o1] {SFT model};
\node[label, below=0.15cm of o2] {Reward model};
\node[label, below=0.15cm of o3] {Aligned LLM};

\end{tikzpicture}

\begin{tikzpicture}[
  node distance=1.0cm and 0.9cm,
  phase/.style={rectangle, draw=hitblue, fill=hitlightblue, rounded corners=4pt,
                minimum width=2.8cm, minimum height=1.6cm, font=\small, align=center,
                text width=2.6cm},
  data/.style={rectangle, draw=hitgold, fill=hitgold!12, rounded corners=3pt,
               minimum width=2.2cm, minimum height=0.55cm, font=\scriptsize, align=center},
  arr/.style={->, >=Stealth, very thick, hitblue},
  darr/.style={->, >=Stealth, thick, hitgold, dashed},
  label/.style={font=\scriptsize, gray, align=center},
]

% Phase 1: SFT
\node[phase] (sft) {\textbf{Phase 1}\\[4pt]Supervised\\Fine-Tuning\\(SFT)};

% Phase 2: Reward Model
\node[phase, right=1.5cm of sft] (rm) {\textbf{Phase 2}\\[4pt]Reward Model\\Training\\(RM)};

% Phase 3: PPO
\node[phase, right=1.5cm of rm] (ppo) {\textbf{Phase 3}\\[4pt]Reinforcement\\Learning\\(PPO)};

% Data nodes above phases
\node[data, above=0.7cm of sft] (d1) {Demonstration\\data};
\node[data, above=0.7cm of rm] (d2) {Human\\preference pairs};
\node[data, above=0.7cm of ppo] (d3) {Prompts\\(unlabeled)};

% Output nodes below phases
\node[data, below=0.7cm of sft] (o1) {$\pi_\text{SFT}$};
\node[data, below=0.7cm of rm] (o2) {$r_\phi(x,y)$};
\node[data, below=0.7cm of ppo] (o3) {$\pi_\theta$ (RLHF)};

% Arrows: data → phase
\draw[darr] (d1) -- (sft);
\draw[darr] (d2) -- (rm);
\draw[darr] (d3) -- (ppo);

% Arrows: phase → output
\draw[arr] (sft) -- (o1);
\draw[arr] (rm) -- (o2);
\draw[arr] (ppo) -- (o3);

% Arrows: phase to phase
\draw[arr] (sft) -- node[above, label] {initialize} (rm);
\draw[arr] (rm) -- node[above, label] {reward\\signal} (ppo);

% KL penalty annotation
\draw[darr, bend left=30] (o1.south) to
  node[below, font=\scriptsize, gray] {KL penalty $\beta$} (o3.south);

% Phase labels below diagram
\node[label, below=0.15cm of o1] {SFT model};
\node[label, below=0.15cm of o2] {Reward model};
\node[label, below=0.15cm of o3] {Aligned LLM};

\end{tikzpicture}

\end{frame}

\begin{frame}{Reward Model Training}
  Given preference pair $(x, \yw, \yl)$ (human says $\yw$ is better):
  \vspace{0.3em}
  \begin{tcolorbox}[hitbox, title=Bradley-Terry Model]
    \[
      P(\yw \succ \yl \given x)
      = \sigma\!\left(r_\phi(x, \yw) - r_\phi(x, \yl)\right)
    \]
    \[
      \mathcal{L}_\text{RM}(\phi) = -\E_{(x,\yw,\yl)}\left[
        \log \sigma\!\left(r_\phi(x, \yw) - r_\phi(x, \yl)\right)
      \right]
    \]
  \end{tcolorbox}
  \vspace{0.3em}
  \begin{itemize}
    \item Reward model $r_\phi$ is usually initialized from the SFT model + a scalar head.
    \item Quality depends heavily on \textbf{annotation quality and diversity}.
    \item Vulnerabilities: reward hacking, out-of-distribution prompts.
  \end{itemize}
  \note{Reward hacking: the RL policy finds ways to score high on $r_\phi$ that don't correspond to genuine quality. This is a fundamental problem.}
\end{frame}

\begin{frame}{PPO Fine-tuning with KL Penalty}
  Maximize expected reward while staying close to the reference SFT model:
  \begin{tcolorbox}[hitbox, title=RLHF Objective]
    \[
      \max_\theta\ \E_{x \sim \mathcal{D},\, y \sim \pitheta(\cdot|x)}\!\left[
        r_\phi(x, y) - \beta \log \frac{\pitheta(y|x)}{\piref(y|x)}
      \right]
    \]
  \end{tcolorbox}
  \vspace{0.3em}
  \begin{itemize}
    \item $\beta > 0$: KL penalty coefficient — controls how far the policy drifts from $\piref$.
    \item Large $\beta$ → safe but conservative; small $\beta$ → risky reward hacking.
    \item PPO (Schulman et al.\ 2017): proximal policy optimization with clipped surrogate.
  \end{itemize}
  \note{The KL penalty is a regularizer. Without it, the policy collapses to repeated reward-hacking outputs.}
\end{frame}

\section{Direct Preference Optimization (DPO)}
\sectionframe{Direct Preference Optimization}{A simpler, stable alternative to RLHF}

\begin{frame}{The DPO Insight}
  \textbf{Rafailov et al.\ (2023)}~\cite{rafailov2023direct}
  \vspace{0.4em}
  \begin{tcolorbox}[hitbox, title=Key Insight]
    The optimal RLHF policy can be written as:
    \[
      \pitheta(y|x) \propto \piref(y|x)\exp\!\left(\tfrac{1}{\beta}r(x,y)\right)
    \]
    Rearranging gives $r(x,y)$ in terms of $\pitheta$ and $\piref$.
    Substituting into the Bradley-Terry model yields a \textbf{closed-form loss}:
  \end{tcolorbox}
  \vspace{0.3em}
  \begin{tcolorbox}[hitbox, title=DPO Loss]
    \[
      \DPOloss(\theta) = -\E\!\left[
        \log \sigma\!\left(
          \beta\log\frac{\pitheta(\yw|x)}{\piref(\yw|x)}
          - \beta\log\frac{\pitheta(\yl|x)}{\piref(\yl|x)}
        \right)
      \right]
    \]
  \end{tcolorbox}
\end{frame}

\begin{frame}{DPO: Advantages and Limitations}
  \begin{columns}[T]
    \column{0.52\textwidth}
      \textbf{Advantages over RLHF:}
      \begin{itemize}
        \item No explicit reward model needed.
        \item No PPO — stable supervised training.
        \item Simpler implementation (classification-style loss).
        \item Competitive or better performance.
      \end{itemize}
    \column{0.44\textwidth}
      \textbf{Limitations:}
      \begin{itemize}
        \item Requires on-policy data or assumption that $\piref \approx \pitheta$.
        \item Cannot update in an online fashion (by default).
        \item Extensions: IPO, KTO, ORPO, SimPO address specific shortcomings.
      \end{itemize}
  \end{columns}
  \vspace{0.5em}
  \textbf{Current state:} DPO (or its variants) is the dominant alignment technique in open-source LLMs.
\end{frame}

\section{Paper Discussions}
\begin{frame}{Paper Discussion 1}
  \papercite{Training Language Models to Follow Instructions with Human Feedback}
    {Ouyang, Wu, Jiang et al.\ (OpenAI)}
    {NeurIPS 2022 — the InstructGPT paper~\cite{ouyang2022training}}
  \vspace{0.4em}
  \textbf{Key findings:}
  \begin{itemize}
    \item 1.3B InstructGPT significantly preferred over 175B GPT-3 by human raters.
    \item Small dataset of human demonstrations + preference labels is sufficient.
    \item ``Alignment tax'': slight performance drop on NLP benchmarks after RLHF.
  \end{itemize}
\end{frame}

\begin{frame}{Paper Discussion 2}
  \papercite{Direct Preference Optimization: Your Language Model is Secretly a Reward Model}
    {Rafailov, Sharma, Mitchell, Ermon, Manning, Finn (Stanford)}
    {NeurIPS 2023~\cite{rafailov2023direct}}
  \vspace{0.4em}
  \textbf{Key insight:} The reward function is implicit in the ratio $\pitheta / \piref$.
  No RL required.\\[0.4em]
  \textbf{Discussion:} If DPO is simpler and competitive, why would anyone still use RLHF?
  What are the scenarios where an explicit reward model is necessary?
\end{frame}

\section{Lab / Demo}
\begin{frame}{Lab Idea: Fine-tuning with SFT and DPO}
  \begin{tcolorbox}[hitbox, title=Lab 3 — Alignment Pipeline]
    \textbf{Goal:} Observe the effect of SFT and DPO on model behavior.\\[4pt]
    \textbf{Tools:} HuggingFace \texttt{trl} library, \texttt{Llama-3.2-1B}
    \begin{enumerate}
      \item Run the base model on 10 prompts (rate helpfulness 1–5).
      \item Apply SFT on a small instruction dataset (e.g., Alpaca 1K).
      \item Apply DPO on a small preference dataset (e.g., Anthropic HH-RLHF 1K).
      \item Re-evaluate on the same prompts. Compare before/after.
    \end{enumerate}
    \textbf{Deliverable:} Table of responses + qualitative analysis.
  \end{tcolorbox}
\end{frame}

\section{Discussion \& Summary}
\begin{frame}{Discussion Questions}
  \begin{enumerate}
    \item RLHF relies on human preference labels. What biases do annotators introduce?
          How could you detect and mitigate this?
    \item The KL penalty in RLHF prevents the model from deviating too far from the reference.
          Is this a safety feature or a limitation? Under what conditions would you increase $\beta$?
    \item DPO optimizes for offline preference data. What happens when the learned policy is much
          better than the reference model — does the DPO loss still provide good gradients?
    \item InstructGPT showed an ``alignment tax.'' Is safety and helpfulness fundamentally at odds
          with raw task performance?
  \end{enumerate}
\end{frame}

\begin{frame}{Summary}
  \begin{itemize}
    \item \textbf{CLM} (next-token prediction) is the dominant pretraining objective.
    \item \textbf{Data quality} matters as much as data quantity; curation is a research area.
    \item \textbf{SFT/Instruction tuning} teaches models to follow instructions.
    \item \textbf{RLHF}: collect preferences → train reward model → PPO with KL penalty.
    \item \textbf{DPO}: derive closed-form loss from RLHF solution; no explicit RL needed.
    \item Modern aligned models (GPT-4, Claude, Llama) use combinations of these techniques.
  \end{itemize}
  \vspace{0.3em}
  \textbf{Next week:} Extending LLMs beyond text — multimodal foundation models.
\end{frame}

\begin{frame}[allowframebreaks]{Suggested Reading}
  \nocite{ouyang2022training, rafailov2023direct, wei2022finetuned,
          brown2020language, touvron2023llama2}
  \bibliography{../../references}
\end{frame}

\end{document}
